% Standalone source-record documentation; no research attempts.
\documentclass[11pt]{article}
\usepackage[margin=25mm]{geometry}
\usepackage{fontspec,unicode-math}
\setmainfont{Latin Modern Roman}
\setmathfont{Latin Modern Math}
\usepackage{amsmath,amssymb,mathtools,xurl,hyperref}
\hypersetup{colorlinks=true,urlcolor=blue,linkcolor=blue}
\setcounter{secnumdepth}{0}
\setlength{\parindent}{0pt}
\setlength{\parskip}{5pt}
\setlength{\emergencystretch}{3em}
\newcommand{\sref}[2]{\hyperlink{source:#1:#2}{[S#2]}}
\newcommand{\cataloguescope}[1]{\paragraph{Source record.}#1}
\newcommand{\floor}[1]{\left\lfloor#1\right\rfloor}
\begin{document}
\section{Causal-Attention Retrieval Motivation}\label{statement}

\subsection{Definitions and mathematical statement}

\paragraph{Shared notation.}Standard number systems are $\mathbb N,\mathbb Z,\mathbb Q,\mathbb R,\mathbb C$. Specialized notation and index ranges follow the cited source and the definitions below.


Causal attention permits a token to depend only on its predecessors under the stated order. Angles on $\mathbb T=\mathbb R/(2\pi\mathbb Z)$ represent circular states. Retrieval requires a specified decoding map linking a later token state to an earlier input token.
\paragraph{Record scope and qualifications.}
The catalogue title is retained as metadata. Named certificates, contracts, bounds or frameworks in that title are not asserted as results of the original source extract.
The supplied registry flags this extraction as unrecoverable. It is displayed as a literal source fragment; the definitions provide context rather than a replacement formal statement.

Its differential equation and decoding convention are incomplete or damaged; the full transformer task is not replaced by a reconstructed simplified theorem.

Source-record limitation: Source flags the statement as unrecoverable: The official report uses this question to define its retrieval observable and then derives an asymptotic formula and a rigorous U-shaped retrieval theorem for the stated model; it is not left as a standalone open problem. No separate clean statement is promoted; the original record remains publicly visible and retained for provenance.
The source text below is rendered literally where necessary; damaged formulas and composite extracts are not silently repaired into a new mathematical claim.
The supplied catalogue marks this record as unrecoverable. Its extract and terminology are documented, but a complete original problem has not been reconstructed.
\paragraph{Statement recorded in the supplied source.}
For a distinguished site i∗ = ⌊\textbackslash{}sigma0 N ⌋, with \textbackslash{}sigma0 \textbackslash{}in (0, 1), we ask whether the last token \textbackslash{}thetaN (t) at time t decodes to the initial token at position i∗ . Large language models based on causal transformers are routinely deployed in regimes where the available context is much longer than the part of the prompt that actually matters for the final prediction. A striking empirical observation is the so-called lost-in-the-middle effect: retrieval is typically best when the rel- evant information is placed near the beginning or the end of the context, and substantially worse when it is placed in the middle [6]. The aim of this work is to give a rigorous mathematical explanation of this phenomenon in a minimal continuous-time model for causal self-attention with recency bias. Our starting point is the interacting particle system viewpoint on transformers developed in [3, 4, 5]. We consider a decoder-type dynamics on the circle T = R/2\textbackslash{}piZ, where the jth token is represented by an angle \textbackslash{}thetaj (t) and only attends to its predecessors. After reducing the original spherical model to dimension 2, and replacing the random softmax denominator by a deterministic normalization, we arrive at j X \textbackslash{}lambda e− N (j−k) \textbackslash{}thetȧj (t) = Pj \textbackslash{}lambda W\textbackslash{}beta′ (\textbackslash{}thetaj (t) − \textbackslash{}thetak (t)), W\textbackslash{}beta (x) = e\textbackslash{}beta cos x . e −N (j−ℓ) k=1 ℓ=1 \sref{statement}{2}
\paragraph{Source status.}
The supplied catalogue labels this source record partially solved; that metadata was checked against the supplied snapshot on 2026-10-11. The label is preserved for provenance and is not an independent verification of present open status or mathematical correctness. The catalogue literature assessment is dated 2026-08-22.
\subsection{Short English statement}
The fragment motivates middle-position retrieval in a continuous attention model.
\subsection{Sources}
\begin{itemize}
\item[\hypertarget{source:statement:1}{S1}] Ulam AI and the UnsolvedMath Contributors, UnsolvedMath: A Curated Collection of Open Mathematics Problems, catalogue entry. Catalogue curation and metadata: CC BY 4.0. The source extract is retained; definitions, English summaries and qualifications are editorial additions.
\url{https://huggingface.co/datasets/ulamai/UnsolvedMath}.
\item[\hypertarget{source:statement:2}{S2}] Flows on Measure Spaces and Applications in Machine Learning (2026), Oberwolfach Reports, DOI 10.4171/owr/2026/13, pp. 795--884
\url{https://doi.org/10.4171/owr/2026/13}.
\item[\hypertarget{source:statement:3}{S3}] Duerinckx--Geshkovski--Rossi, Kinetic theory for Transformers
\url{https://arxiv.org/abs/2605.09213}.
\item[\hypertarget{source:statement:4}{S4}] Liu et al., Lost in the Middle
\url{https://doi.org/10.1162/tacl_a_00638}.
\item Source-specific Oberwolfach Reports licences, including CC BY-SA 4.0 where applicable, are retained; see THIRD\_PARTY\_NOTICES.md.
\end{itemize}
\paragraph{Preparation and review.}
The original source and UnsolvedMath contributors retain attribution for the source extract. Codex (GPT-6; exact serving revision not exposed) assisted with definitions, English summaries, scope notes and formatting. No new human mathematical authorship or independently verified solution is asserted. Catalogue identity checked against the complete supplied registry and current repository. Distinct upstream records, including repeated questions and broad or damaged extracts, are retained for complete coverage of the source CS category. A repeated record is not a new independent conjecture. Complete mathematical-equivalence and current-literature review has not been performed. Preparation and final review follow contribution prompts C and D, with source gaps disclosed. Structural and rendering checks do not establish a new theorem or open status.
\label{end:statement}
\end{document}
